\documentclass[11pt,a4paper]{article}

% ============================================================
% Guía de instalación y configuración del entorno de trabajo
% Programación y Base de Datos - Proyecto integrador FIRELAB-Loja
% ============================================================

\usepackage[utf8]{inputenc}
\usepackage[T1]{fontenc}
\usepackage[spanish]{babel}
\usepackage{lmodern}
\usepackage[margin=2.3cm]{geometry}
\usepackage{amsmath,amssymb}
\usepackage{graphicx}
\usepackage{xcolor}
\usepackage{enumitem}
\usepackage[most]{tcolorbox}
\usepackage{fancyhdr}
\usepackage{titlesec}
\usepackage{tocloft}
\setlength{\cftsecnumwidth}{2.3em}
\usepackage{hyperref}
\usepackage{parskip}

% ---------------- Paleta (misma identidad que Programación) ----------------
\definecolor{morado}{RGB}{74,20,99}
\definecolor{moradoOsc}{RGB}{48,11,68}
\definecolor{moradoMedio}{RGB}{123,63,157}
\definecolor{moradoClaro}{RGB}{224,208,235}
\definecolor{tealProfundo}{RGB}{0,90,99}
\definecolor{tealClaro}{RGB}{204,232,232}
\definecolor{acento}{RGB}{198,110,38}
\definecolor{acentoClaro}{RGB}{247,226,201}
\definecolor{alerta}{RGB}{165,40,40}
\definecolor{alertaClara}{RGB}{242,213,208}
\definecolor{grafito}{RGB}{51,51,51}
\definecolor{grisClaro}{RGB}{244,246,245}

\hypersetup{
  colorlinks=true,
  linkcolor=morado,
  urlcolor=moradoMedio,
  citecolor=morado,
  pdftitle={Guía de instalación y configuración del entorno de trabajo},
  pdfauthor={Programación y Base de Datos},
  bookmarksopen=true
}

\titleformat{\section}{\Large\bfseries\color{morado}}{\thesection.}{0.5em}{}
\titleformat{\subsection}{\large\bfseries\color{moradoOsc}}{\thesubsection}{0.5em}{}
\titlespacing*{\section}{0pt}{1.4em}{0.8em}
\titlespacing*{\subsection}{0pt}{1.1em}{0.5em}

\pagestyle{fancy}
\fancyhf{}
\renewcommand{\headrulewidth}{0.6pt}
\renewcommand{\footrulewidth}{0.4pt}
\fancyhead[L]{\small\color{grafito}Programación y Base de Datos -- FIRELAB-Loja}
\fancyhead[R]{\small\color{grafito}Guía de instalación}
\fancyfoot[C]{\small\thepage}

\setlist[itemize]{itemsep=2pt, topsep=3pt}
\setlist[enumerate]{itemsep=4pt, topsep=3pt}

% ---------------- Cajas de aviso ----------------
\newtcolorbox{tip}[1][]{colback=tealClaro!45,colframe=tealProfundo,arc=2pt,boxrule=0.9pt,
  left=8pt,right=8pt,top=6pt,bottom=6pt,title=Consejo,fonttitle=\bfseries,coltitle=white,
  colbacktitle=tealProfundo,#1}
\newtcolorbox{warn}[1][]{colback=acentoClaro!70,colframe=acento,arc=2pt,boxrule=0.9pt,
  left=8pt,right=8pt,top=6pt,bottom=6pt,title=Atención,fonttitle=\bfseries,coltitle=white,
  colbacktitle=acento,#1}
\newtcolorbox{crit}[1][]{colback=alertaClara!60,colframe=alerta,arc=2pt,boxrule=1.1pt,
  left=8pt,right=8pt,top=6pt,bottom=6pt,title=¡Muy importante!,fonttitle=\bfseries,coltitle=white,
  colbacktitle=alerta,#1}
\newtcolorbox{verify}[1][]{colback=moradoClaro!35,colframe=morado,arc=2pt,boxrule=0.9pt,
  left=8pt,right=8pt,top=6pt,bottom=6pt,title=Verifique antes de continuar,fonttitle=\bfseries,coltitle=white,
  colbacktitle=morado,#1}
\newtcolorbox{stepbox}[1][]{colback=grisClaro,colframe=moradoMedio,arc=2pt,boxrule=0.7pt,
  left=8pt,right=8pt,top=6pt,bottom=6pt,#1}

\newcommand{\code}[1]{\texttt{#1}}
\newcommand{\link}[1]{\url{#1}}

\begin{document}
\shorthandoff{"}

% ============================================================
% Portada
% ============================================================
\begin{titlepage}
\centering
\vspace*{1.5cm}
{\Huge\bfseries\color{morado} Guía de instalación\\[4pt] y configuración del entorno}\\[1.2cm]
{\Large\color{moradoOsc} Programación y Base de Datos -- Unidad 1}\\[6pt]
{\large\color{grafito} Proyecto integrador FIRELAB-Loja}\\[2cm]

\begin{tcolorbox}[colback=moradoClaro!30,colframe=morado,arc=3pt,boxrule=1pt,width=0.82\textwidth,
  left=14pt,right=14pt,top=10pt,bottom=10pt]
\textbf{¿Qué instalará con esta guía?}
\begin{itemize}[itemsep=3pt]
  \item Visual Studio Code (editor de código)
  \item Python y las librerías del curso (\code{pandas}, \code{matplotlib})
  \item Git (control de versiones) y una cuenta de GitHub
  \item El repositorio de trabajo de su equipo para el proyecto FIRELAB-Loja
\end{itemize}
\end{tcolorbox}

\vspace{1.3cm}
{\small\color{grafito} Tiempo estimado: 45--60 minutos \quad\textbullet\quad Hágalo con conexión a internet estable}\\[0.4cm]
{\small\color{grafito} Traiga su equipo con esta instalación completa a la primera sesión práctica}

\vfill
{\small\color{grafito} Universidad Nacional de Loja -- Carrera de Ingeniería Ambiental}
\end{titlepage}

\tableofcontents
\newpage

% ============================================================
\section{Antes de empezar}
% ============================================================

Esta guía lo llevará, paso a paso, desde un computador recién formateado hasta tener listo el entorno completo de trabajo del curso: editor, lenguaje, librerías, control de versiones y el repositorio de su equipo para el proyecto integrador.

\begin{tip}
Siga los pasos \textbf{en el orden en que aparecen}. Cada paso depende del anterior (por ejemplo, VS Code necesita que Python ya esté instalado para poder usarlo). Si se salta un paso, es probable que algo falle más adelante.
\end{tip}

\subsection*{Requisitos previos}
\begin{itemize}
  \item Un computador con Windows 10/11, macOS o Linux, con al menos 3\,GB de espacio libre.
  \item Conexión a internet estable (las descargas suman cerca de 300--400\,MB en total).
  \item Permisos de administrador en el equipo (para instalar programas).
  \item Una dirección de correo electrónico personal (para crear la cuenta de GitHub).
\end{itemize}

\begin{crit}
Todos los pasos marcados como \textbf{¡Atención en Windows!} son la causa más frecuente de que "nada funcione" en la primera sesión práctica. Léalos con calma, no los salte.
\end{crit}

% ============================================================
\section{Instalar Visual Studio Code}
% ============================================================

Visual Studio Code (VS Code) es el editor de código que usaremos durante todo el curso.

\begin{stepbox}
\textbf{Paso 1.} Abra su navegador y vaya a:
\begin{center}\link{https://code.visualstudio.com/download}\end{center}
La página detecta automáticamente su sistema operativo y le muestra el botón de descarga correcto (Windows, macOS o Linux). Haga clic en ese botón.
\end{stepbox}

\begin{stepbox}
\textbf{Paso 2.} Ejecute el instalador descargado.
\begin{itemize}
  \item \textbf{Windows:} abra el archivo \code{VSCodeUserSetup-*.exe}. Acepte el acuerdo de licencia y deje todas las casillas marcadas por defecto en la pantalla de "Tareas adicionales" (incluida \textbf{"Add to PATH"}, si aparece). Pulse "Siguiente" hasta finalizar.
  \item \textbf{macOS:} abra el archivo \code{.zip} descargado, arrastre \code{Visual Studio Code.app} a la carpeta \code{Aplicaciones}.
  \item \textbf{Linux (Ubuntu/Debian):} instale el paquete \code{.deb} descargado con \code{sudo apt install ./nombre-del-archivo.deb}, o siga las instrucciones específicas de su distribución en la misma página de descarga.
\end{itemize}
\end{stepbox}

\begin{verify}
Abra Visual Studio Code desde el menú de inicio (Windows), Launchpad (macOS) o su lanzador de aplicaciones (Linux). Debe abrirse una ventana de bienvenida de VS Code. Si esto ocurre, este paso está completo.
\end{verify}

% ============================================================
\section{Instalar Python}
% ============================================================

Python es el lenguaje de programación que usaremos en todo el curso.

\begin{stepbox}
\textbf{Paso 3.} Vaya a:
\begin{center}\link{https://www.python.org/downloads/}\end{center}
Haga clic en el botón amarillo que descarga la última versión estable (por ejemplo, "Download Python 3.12.x"). Evite versiones marcadas como "pre-release" o "rc".
\end{stepbox}

\begin{crit}
\textbf{¡Atención en Windows!} En la \textbf{primera pantalla} del instalador aparecen dos casillas en la parte inferior: \textbf{"Add python.exe to PATH"} y "Install launcher for all users". \textbf{Marque "Add python.exe to PATH" antes} de hacer clic en "Install Now".

Si no marca esa casilla, Windows no sabrá dónde encontrar Python, y VS Code tampoco. Los síntomas típicos son mensajes como \code{'python' no se reconoce como un comando interno o externo}. Si esto le pasa, desinstale Python desde "Agregar o quitar programas" y vuelva a instalarlo marcando la casilla correctamente.
\end{crit}

\begin{stepbox}
\textbf{macOS:} el instalador \code{.pkg} no tiene la casilla de PATH (macOS lo configura automáticamente); simplemente siga los pasos del asistente con las opciones por defecto.

\textbf{Linux:} la mayoría de distribuciones ya incluyen Python 3. Verifique con el comando del Paso 4; si falta, instálelo con el gestor de paquetes de su distribución (por ejemplo \code{sudo apt install python3 python3-pip} en Ubuntu/Debian).
\end{stepbox}

\begin{verify}
\textbf{Paso 4.} Abra una terminal:
\begin{itemize}
  \item \textbf{Windows:} busque "PowerShell" o "Símbolo del sistema" en el menú de inicio.
  \item \textbf{macOS/Linux:} busque "Terminal" en el lanzador de aplicaciones.
\end{itemize}
Escriba (en Windows) \code{python --version} o (en macOS/Linux) \code{python3 --version} y presione Enter. Debe aparecer algo como:
\begin{center}\code{Python 3.12.4}\end{center}
Si en cambio aparece un error de "comando no reconocido", revise el aviso de PATH de arriba.
\end{verify}

% ============================================================
\section{Configurar VS Code para trabajar con Python}
% ============================================================

\begin{stepbox}
\textbf{Paso 5.} Abra Visual Studio Code. En la barra lateral izquierda, haga clic en el ícono de extensiones (cuatro cuadrados, el cuarto separado) o use el atajo \code{Ctrl+Shift+X} (Windows/Linux) o \code{Cmd+Shift+X} (macOS).
\end{stepbox}

\begin{stepbox}
\textbf{Paso 6.} Busque e instale (botón azul "Install" junto a cada resultado) estas dos extensiones oficiales de Microsoft:
\begin{itemize}
  \item \textbf{Python} -- \link{https://marketplace.visualstudio.com/items?itemName=ms-python.python}
  \item \textbf{Jupyter} -- \link{https://marketplace.visualstudio.com/items?itemName=ms-toolsai.jupyter}
\end{itemize}
Busque exactamente por el nombre \code{Python} y \code{Jupyter}; ambas son publicadas por "Microsoft" (verifique el editor/publisher en la ficha de la extensión antes de instalar).
\end{stepbox}

\begin{stepbox}
\textbf{Paso 7.} Cree un archivo de prueba:
\begin{enumerate}
  \item En VS Code, vaya a \textbf{Archivo $\to$ Nuevo archivo de texto}.
  \item Guárdelo (\code{Ctrl+S}) con el nombre \code{prueba.py} en una carpeta de su elección (por ejemplo, cree una carpeta \code{Escritorio/prueba\_python}).
  \item Escriba esta línea: \code{print("Hola FIRELAB")}
  \item En la esquina inferior derecha de VS Code debe aparecer la versión de Python detectada (por ejemplo, "3.12.4"). Si dice "Select Interpreter", haga clic ahí y elija la versión de Python instalada en el Paso 3.
  \item Ejecute el archivo con el botón $\blacktriangleright$ (arriba a la derecha) o clic derecho $\to$ "Run Python File in Terminal".
\end{enumerate}
\end{stepbox}

\begin{verify}
En la terminal integrada, en la parte inferior de VS Code, debe aparecer el texto \code{Hola FIRELAB}. Si lo ve, VS Code y Python ya están correctamente conectados.
\end{verify}

% ============================================================
\section{Instalar las librerías del curso}
% ============================================================

\begin{stepbox}
\textbf{Paso 8.} Abra la terminal integrada de VS Code: menú \textbf{Terminal $\to$ Nueva terminal}, o el atajo \code{Ctrl+\textasciigrave} (comilla invertida, tecla arriba del Tab).

Escriba el siguiente comando y presione Enter:
\begin{center}\code{pip install pandas matplotlib}\end{center}
(En macOS/Linux, si \code{pip} no funciona, use \code{pip3 install pandas matplotlib}.)

Espere a que termine la instalación; puede tardar uno o dos minutos.
\end{stepbox}

\begin{tip}
\code{sqlite3} \textbf{no se instala}: viene incluido de fábrica con Python, listo para usar con \code{import sqlite3}.
\end{tip}

\begin{verify}
\textbf{Paso 9.} En su archivo \code{prueba.py} (o uno nuevo, \code{prueba\_final.py}), reemplace el contenido por:
\begin{center}\code{import pandas; import matplotlib; import sqlite3; print("Todo listo")}\end{center}
Ejecútelo. Si aparece \code{Todo listo} sin ningún error en rojo, la instalación de Python y sus librerías está completa.
\end{verify}

% ============================================================
\section{Crear una cuenta de GitHub}
% ============================================================

GitHub es la plataforma donde su equipo alojará el código y los datos del proyecto integrador, y donde el docente revisará su avance.

\begin{stepbox}
\textbf{Paso 10.} Vaya a:
\begin{center}\link{https://github.com/signup}\end{center}
Regístrese con su correo electrónico personal (no institucional, para poder seguir usando la cuenta después de graduarse). Elija un \textbf{nombre de usuario profesional}: evite apodos; use algo basado en su nombre real (por ejemplo, \code{cguillermo-chuncho} en vez de \code{cracker123}), ya que este perfil puede formar parte de su portafolio profesional a futuro.
\end{stepbox}

\begin{tip}
GitHub ofrece cuentas gratuitas con repositorios privados ilimitados para estudiantes y en general. No es necesario pagar nada para este curso. Adicionalmente, con un correo institucional (.edu) puede solicitar el \href{https://education.github.com/pack}{GitHub Student Developer Pack}, que da acceso a herramientas profesionales gratuitas.
\end{tip}

\begin{verify}
Inicie sesión en \link{https://github.com} con su nuevo usuario y contraseña. Debe ver su panel principal ("Dashboard") de GitHub.
\end{verify}

% ============================================================
\section{Instalar Git}
% ============================================================

Git es el programa que registra el historial de cambios de su código y lo sincroniza con GitHub. GitHub es la plataforma web; Git es la herramienta que se instala en su computador.

\begin{stepbox}
\textbf{Paso 11.} Descargue Git desde:
\begin{center}\link{https://git-scm.com/downloads}\end{center}
Elija su sistema operativo y descargue el instalador.
\begin{itemize}
  \item \textbf{Windows:} ejecute el instalador y acepte las opciones por defecto en todas las pantallas (son más de diez pantallas; para este curso los valores predeterminados funcionan bien). Un detalle a revisar: en la pantalla "Adjusting your PATH environment", deje seleccionada la opción recomendada ("Git from the command line and also from 3rd-party software").
  \item \textbf{macOS:} si tiene Xcode Command Line Tools, es posible que Git ya esté instalado. Si no, el instalador de la página oficial funciona igual de bien.
  \item \textbf{Linux:} instale con \code{sudo apt install git} (Ubuntu/Debian) o el gestor de paquetes equivalente.
\end{itemize}
\end{stepbox}

\begin{verify}
\textbf{Paso 12.} Abra una terminal nueva (ciérrela y ábrala de nuevo si ya tenía una abierta) y escriba:
\begin{center}\code{git --version}\end{center}
Debe mostrar algo como \code{git version 2.45.0}.
\end{verify}

\begin{stepbox}
\textbf{Paso 13.} Configure su identidad en Git (una sola vez, en este computador). En la terminal, ejecute estas dos líneas reemplazando sus datos:
\begin{center}
\code{git config --global user.name "Su Nombre Completo"}\\
\code{git config --global user.email "su\_correo@ejemplo.com"}
\end{center}
Use \textbf{el mismo correo} con el que se registró en GitHub, para que sus contribuciones queden asociadas a su perfil.
\end{stepbox}

% ============================================================
\section{Conectar VS Code con GitHub}
% ============================================================

\begin{stepbox}
\textbf{Paso 14.} En VS Code, haga clic en el ícono de "Control de código fuente" (Source Control) en la barra lateral izquierda (el ícono con forma de rama, o el atajo \code{Ctrl+Shift+G}).
\end{stepbox}

\begin{stepbox}
\textbf{Paso 15.} La primera vez que intente subir código a GitHub desde VS Code (más adelante, en el Paso 17), aparecerá una ventana pidiéndole iniciar sesión. Haga clic en \textbf{"Sign in with GitHub"} y autorice el acceso en la ventana del navegador que se abre. No necesita configurar nada más de forma manual: VS Code gestiona la autenticación por usted.
\end{stepbox}

\begin{tip}
Si en algún momento le pide usuario y contraseña de Git \textbf{en la terminal} (no en una ventana de VS Code) y su contraseña normal de GitHub no funciona, es porque GitHub ya no acepta contraseñas por línea de comandos: use el inicio de sesión de VS Code (Paso 15) en vez de escribir credenciales manualmente.
\end{tip}

% ============================================================
\section{Construir el repositorio de su equipo}
% ============================================================

Cada pareja o grupo de equipos trabaja un bloque temático del dataset FIRELAB-Loja (cantón Loja, 2019--2024). Use esta tabla para identificar el suyo:

\begin{center}
\renewcommand{\arraystretch}{1.3}
\begin{tabular}{c p{9.5cm}}
\hline
\textbf{Equipos} & \textbf{Bloque temático} \\
\hline
1 y 2 & Relieve y accesibilidad \\
3 & Cobertura vegetal y uso del suelo \\
4 y 5 & Clima y condiciones atmosféricas \\
6 & Incendios forestales y riesgo \\
\hline
\end{tabular}
\end{center}

\begin{stepbox}
\textbf{Paso 16 -- Crear el repositorio en GitHub.}
\begin{enumerate}
  \item Inicie sesión en \link{https://github.com} y haga clic en el botón verde \textbf{"New"} (o el signo \code{+} en la esquina superior derecha $\to$ "New repository").
  \item \textbf{Nombre del repositorio:} use el patrón \code{firelab-loja-equipoN-bloque}, reemplazando \code{N} por su número de equipo y \code{bloque} por una palabra corta de su tema. Ejemplos: \code{firelab-loja-equipo1-relieve}, \code{firelab-loja-equipo3-cobertura}, \code{firelab-loja-equipo6-incendios}.
  \item \textbf{Visibilidad:} elija \textbf{Private} (privado), y agregue como colaboradores a sus compañeros de equipo y al docente (Settings $\to$ Collaborators, una vez creado).
  \item Marque la casilla \textbf{"Add a README file"}.
  \item Haga clic en \textbf{"Create repository"}.
\end{enumerate}
\end{stepbox}

\begin{stepbox}
\textbf{Paso 17 -- Clonar el repositorio en su computador con VS Code.}
\begin{enumerate}
  \item En VS Code, presione \code{Ctrl+Shift+P} (o \code{Cmd+Shift+P} en macOS) para abrir la paleta de comandos.
  \item Escriba \code{Git: Clone} y selecciónelo.
  \item Elija \textbf{"Clone from GitHub"}, inicie sesión si se lo pide (vea el Paso 15), y busque su repositorio \code{firelab-loja-equipoN-bloque} en la lista.
  \item Elija una carpeta local donde guardarlo (por ejemplo, \code{Documentos/proyectos}).
  \item Cuando termine, VS Code preguntará si desea abrir el repositorio clonado: elija \textbf{"Open"}.
\end{enumerate}
\end{stepbox}

\begin{stepbox}
\textbf{Paso 18 -- Estructura de carpetas recomendada.} Dentro de su repositorio recién clonado, cree esta estructura (clic derecho en el panel de "Explorador" de VS Code $\to$ "New Folder" / "New File"):
\begin{center}
\begin{minipage}{0.75\textwidth}
\ttfamily
firelab-loja-equipoN-bloque/\\
\hspace*{1em}datos/\quad\rmfamily\itshape (datos crudos y procesados de su bloque)\\
\ttfamily\hspace*{1em}notebooks/\quad\rmfamily\itshape (exploración en Jupyter, archivos .ipynb)\\
\ttfamily\hspace*{1em}src/\quad\rmfamily\itshape (funciones y scripts .py reutilizables)\\
\ttfamily\hspace*{1em}resultados/\quad\rmfamily\itshape (gráficos y tablas que generen)\\
\ttfamily\hspace*{1em}README.md
\end{minipage}
\end{center}
\end{stepbox}

\begin{stepbox}
\textbf{Paso 19 -- Completar el README y su primer commit.} Abra el archivo \code{README.md} y escriba una descripción breve: nombre del equipo, integrantes, bloque temático asignado y la pregunta guía del proyecto (tómela de la diapositiva "Proyecto integrador" de la presentación del curso). Guarde el archivo.

Luego, en el panel de \textbf{Control de código fuente} (Paso 14):
\begin{enumerate}
  \item Escriba un mensaje corto en el recuadro de arriba, por ejemplo: \code{Estructura inicial del proyecto y README}.
  \item Haga clic en el botón \textbf{"Commit"} (o el ícono de check $\checkmark$).
  \item Haga clic en \textbf{"Sync Changes"} (o "Push") para subir los cambios a GitHub.
\end{enumerate}
\end{stepbox}

\begin{verify}
Abra su repositorio en \link{https://github.com} desde el navegador y actualice la página. Debe ver la carpeta \code{datos/}, \code{notebooks/}, \code{src/}, \code{resultados/} y su \code{README.md} ya con el contenido que escribió. Si es así, su equipo ya tiene el proyecto listo para empezar a trabajar.
\end{verify}

% ============================================================
\section{Lista de verificación final}
% ============================================================

Antes de la próxima sesión práctica, confirme que puede marcar \textbf{todos} estos puntos:

\begin{itemize}[label=$\square$]
  \item Visual Studio Code instalado y abre correctamente.
  \item \code{python --version} (o \code{python3 --version}) muestra una versión 3.11 o superior.
  \item Las extensiones \textbf{Python} y \textbf{Jupyter} instaladas en VS Code.
  \item Un archivo \code{.py} con \code{print(...)} se ejecuta sin errores desde VS Code.
  \item \code{pip install pandas matplotlib} se ejecutó sin errores.
  \item Cuenta de GitHub creada, con nombre de usuario profesional.
  \item \code{git --version} funciona, y \code{git config --global user.name/user.email} configurados.
  \item Repositorio \code{firelab-loja-equipoN-bloque} creado en GitHub, clonado en su computador, con la estructura de carpetas y el primer commit ya subido ("pusheado").
\end{itemize}

% ============================================================
\section{Problemas frecuentes}
% ============================================================

\begin{description}
  \item[\textcolor{alerta}{\textbf{"python" no se reconoce como un comando}}] Python se instaló sin marcar "Add python.exe to PATH". Desinstale Python desde el Panel de control y reinstálelo marcando esa casilla (Sección 3).

  \item[\textcolor{alerta}{\textbf{"pip" no se reconoce}}] Mismo origen que el anterior: revise el PATH de Python. Como alternativa temporal, use \code{python -m pip install pandas matplotlib}.

  \item[\textcolor{alerta}{\textbf{VS Code no detecta el intérprete de Python}}] Con un archivo \code{.py} abierto, haga clic en la esquina inferior derecha (donde dice la versión de Python o "Select Interpreter") y elija manualmente la versión instalada.

  \item[\textcolor{alerta}{\textbf{"git" no se reconoce como un comando}}] Cierre todas las ventanas de terminal y VS Code, y vuelva a abrirlas (Git actualiza el PATH del sistema, pero las terminales ya abiertas no se enteran hasta reiniciarlas).

  \item[\textcolor{alerta}{\textbf{No puedo hacer "push" a GitHub (permiso denegado)}}] Verifique que inició sesión con "Sign in with GitHub" dentro de VS Code (Sección 8), y que su usuario de GitHub fue agregado como colaborador del repositorio si no es el dueño.

  \item[\textcolor{alerta}{\textbf{Instalé todo pero algo más no funciona}}] Revise primero la Lista de verificación final (Sección 10) para aislar cuál paso específico falló, y traiga su equipo con el error visible en pantalla a la próxima sesión práctica.
\end{description}

\vspace{0.6cm}
\begin{tcolorbox}[colback=moradoClaro!30,colframe=morado,arc=3pt,boxrule=1pt]
\textbf{¿Sigue con problemas?} Consulte en el horario de tutoría del curso, o escriba al docente antes de la sesión práctica adjuntando una captura de pantalla del error. Es normal que la primera instalación tenga algún tropiezo -para eso está esta guía y el acompañamiento en clase.
\end{tcolorbox}

\end{document}
